% LaTeX template for the pandoc-rendered reports.
%
% Pandoc does the structural conversion; this file does the typography. It
% is deliberately minimal and self-contained: the packages below are the ones
% pandoc's LaTeX output actually needs, and nothing else, because every extra
% package is another thing that can be missing from a TeX bundle.
%
% The two .tex papers in this project use their own class files; this
% template is for the prose reports, which are markdown first and PDF second.
$if(bodypt12)$
\documentclass[12pt,a4paper]{article}
$else$
\documentclass[11pt,a4paper]{article}
$endif$

\usepackage[T1]{fontenc}
\usepackage[utf8]{inputenc}
$if(timesfont)$
% Times for text and for maths, so a formula in the PDF is set in the same
% face as the paragraph around it. The `lmodern` below is then redundant and
% is not loaded, because the two would fight over the roman face.
\usepackage{newtxtext}
\usepackage{newtxmath}
$else$
\usepackage{lmodern}
$endif$
$if(widepage)$
% A document whose content is a full-width picture of a slide needs the
% page to be wider than the prose margin, or the picture is too small to
% read a table from.
\usepackage[margin=1.6cm,top=2.2cm,bottom=2.1cm]{geometry}
$else$
\usepackage[margin=2.4cm,top=2.6cm,bottom=2.5cm]{geometry}
$endif$
\usepackage{amsmath}
$if(timesfont)$
% amssymb is not loaded alongside newtxmath: both define \Bbbk, and the
% clash is a hard error, not a warning. newtxmath carries the symbols the
% documents use.
$else$
\usepackage{amssymb}
$endif$
\usepackage{array}
\usepackage{booktabs}
\usepackage{longtable}
\usepackage{calc}
\usepackage{graphicx}
\usepackage[table]{xcolor}
\usepackage{framed}
\usepackage{fancyvrb}
\usepackage{listings}
\usepackage{titlesec}
\usepackage{fancyhdr}
\usepackage{enumitem}
\usepackage[hidelinks]{hyperref}

% NOTE: the `caption` package is deliberately NOT loaded. It redefines the
% float machinery in a way that breaks pandoc's longtable wrapper -- the
% tables come out with "No counter 'none' defined" and the document does not
% compile. The reports carry no figures, so nothing is lost by leaving it
% out; if a figure is ever added, it will need a different fix than the
% obvious one.

% -- typography ----------------------------------------------------------
% The theme is monochrome. These documents are printed, photocopied and
% read on a projector in a room where the colour on a screen is not
% guaranteed to survive, and a document whose meaning depends on a hue is
% a document that fails on a mono printer. The hierarchy is carried by
% weight, size and greyscale instead: black for the title and the section
% headings, dark grey for the subsections, mid grey for the rules and the
% running head, and a very light grey for the shading.
\definecolor{brandnavy}{HTML}{000000}
\definecolor{brandaccent}{HTML}{333333}
\definecolor{shadecolor}{RGB}{242,242,242}
\definecolor{rulegray}{HTML}{8C8C8C}

% The section numbers are written into the markdown headings themselves, so
% LaTeX must not add a second one. The text refers to sections by number
% ("see 5.4"), and those numbers have to be the numbers on the page.
\setcounter{secnumdepth}{0}

\titleformat{\section}{\normalfont\Large\bfseries\color{brandnavy}}{\thesection}{0.7em}{}
\titleformat{\subsection}{\normalfont\large\bfseries\color{brandaccent}}{\thesubsection}{0.6em}{}
\titleformat{\subsubsection}{\normalfont\normalsize\bfseries\color{brandnavy}}{\thesubsubsection}{0.6em}{}
\titlespacing*{\section}{0pt}{16pt}{6pt}
\titlespacing*{\subsection}{0pt}{12pt}{4pt}

\setlist{itemsep=2pt,topsep=3pt,parsep=0pt}
\setlength{\LTcapwidth}{\textwidth}
\renewcommand{\arraystretch}{1.15}

% -- code blocks ---------------------------------------------------------
% Pandoc is invoked with --no-highlight, so a code block arrives as a plain
% verbatim environment and no token macros are needed. The definitions
% below are the fallbacks for the case where a document is rendered with
% highlighting switched on.
\providecommand{\passthrough}[1]{#1}

% Pandoc emits \tightlist for a list with no blank lines between its items.
% Its own template defines it; this template has to as well.
\providecommand{\tightlist}{%
  \setlength{\itemsep}{0pt}\setlength{\parskip}{0pt}}

% Pandoc also wraps every image in \pandocbounded, which its own template
% defines. Without it, any document carrying a figure fails to compile with
% "Undefined control sequence" -- and the figures are the point.
\providecommand{\pandocbounded}[1]{#1}
\lstset{basicstyle=\ttfamily\small,breaklines=true,frame=single,rulecolor=\color{rulegray},columns=fullflexible,keepspaces=true,showstringspaces=false,upquote=true}

% -- page furniture ------------------------------------------------------
\pagestyle{fancy}
\fancyhf{}
\renewcommand{\headrulewidth}{0.4pt}
\renewcommand{\footrulewidth}{0.4pt}
% The running head is short on purpose: the full title and the author
% together overrun the width of the page and overlap.
$if(deptfurniture)$
% The department's own page furniture: no running head, and the page number
% bottom right, which is what the reference report has.
\fancyhead[L]{}
\fancyhead[R]{}
\renewcommand{\headrulewidth}{0pt}
\fancyfoot[R]{\small\color{brandaccent}\thepage}
$else$
\fancyhead[L]{\small\color{brandaccent}Evidence-Grounded RAG for Industrial Decision Support}
\fancyhead[R]{\small\color{gray}$if(shorttitle)$$shorttitle$$endif$}
\fancyfoot[C]{\small\color{brandaccent}\thepage}
$endif$
\setlength{\headheight}{14pt}

\begin{document}

$if(frontmatter)$
% ---------------------------------------------------------------- front matter
% The title page and the certificate page, in the department's template. The
% same two pages are built in the Word file; this is the reading copy of them,
% so the two must carry the same text.
\pagestyle{empty}
\begin{center}
\vspace*{0.4cm}
{\fontsize{20}{24}\selectfont\bfseries AN EVIDENCE-GROUNDED\par}
\vspace{0.1cm}
{\fontsize{18}{22}\selectfont\bfseries RETRIEVAL-AUGMENTED GENERATION\par}
\vspace{0.1cm}
{\fontsize{18}{22}\selectfont\bfseries FRAMEWORK\par}
\vspace{0.1cm}
{\fontsize{16}{20}\selectfont\bfseries FOR INDUSTRIAL DECISION SUPPORT\par}
\vspace{0.5cm}
{\fontsize{14}{17}\selectfont\itshape M.Tech Project Report\par}
\vspace{0.6cm}
\includegraphics[width=5.66cm]{assets/iitkgp_emblem.png}\par
\vspace{1.5cm}
{\fontsize{13}{16}\selectfont By\par}
\vspace{0.2cm}
{\fontsize{15}{18}\selectfont\bfseries Ritik Raj [25IM60R12]\par}
\vspace{0.35cm}
{\fontsize{13}{16}\selectfont Under the supervision of\par}
\vspace{0.15cm}
{\fontsize{14}{17}\selectfont\bfseries Dr. Sri Krishna Kumar\par}
\vspace{0.55cm}
{\fontsize{12}{15}\selectfont\bfseries INDIAN INSTITUTE OF TECHNOLOGY KHARAGPUR)\par}
\vspace{0.1cm}
{\fontsize{12}{15}\selectfont\bfseries DEPARTMENT OF INDUSTRIAL AND SYSTEMS ENGINEERING\par}
\vspace{0.1cm}
{\fontsize{12}{15}\selectfont\bfseries INDIAN INSTITUTE OF TECHNOLOGY, KHARAGPUR\par}
\vspace{0.1cm}
{\fontsize{12}{15}\selectfont KHARAGPUR- 721302, INDIA\par}
\vspace{0.3cm}
{\fontsize{12}{15}\selectfont\bfseries NOVEMBER, 2026\par}
\end{center}
\newpage

% The certificate page. The emblem sits in the margin beside the letterhead,
% the way the department's template has it.
\fancypagestyle{deptcert}{\fancyhf{}\renewcommand{\headrulewidth}{0pt}\renewcommand{\footrulewidth}{0pt}\fancyfoot[R]{\small\thepage}}
\thispagestyle{deptcert}
\noindent
\begin{minipage}[t]{2.4cm}
  \includegraphics[width=1.81cm]{assets/iitkgp_emblem.png}
\end{minipage}%
\begin{minipage}[t]{12.2cm}
  \fontsize{11}{13}\selectfont\bfseries
  DEPARTMENT OF INDUSTRIAL AND SYSTEMS ENGINEERING\\
  INDIAN INSTITUTE OF TECHLOGY KHARAGPUR,\\
  KHARAGPUR, INDIA-721302
\end{minipage}
\par\vspace{0.35cm}
\noindent\rule{\textwidth}{0.7pt}
\par\vspace{1.1cm}
\begin{center}
{\fontsize{18}{22}\selectfont\bfseries CERTIFICATE\par}
\end{center}
\vspace{1.1cm}
{\fontsize{12}{20.6}\selectfont
This is to certify that this report entitled An Evidence-Grounded Retrieval-Augmented
Generation Framework for Industrial Decision Support submitted by Ritik Raj [25IM60R12]
to the Indian Institute of Technology Kharagpur is a record of the bona fide project
work carried out under my supervision and is worthy of consideration for the MTP
Project. This project has fulfilled all the requirements as per the regulations of
the Institute and has reached the standard needed for the submission.\par}
\vfill
{\fontsize{10.6}{13}\selectfont
\textbf{Dr. Sri Krishna Kumar}\\
Department of Industrial and Systems Engineering\\
Indian Institute of Technology Kharagpur\\
India-721302\\
Date: 18-11-2026\par}
\newpage
\pagestyle{fancy}
$else$
$if(title)$
\begin{center}
  {\LARGE\bfseries\color{brandnavy} $title$ \par}
  \vspace{0.4em}
  $if(subtitle)${\large\color{brandaccent} $subtitle$ \par}\vspace{0.6em}$endif$
  $for(author)${\large $author$ \par}$endfor$
  $if(date)${\small\color{gray} $date$ \par}$endif$
\end{center}
\vspace{1.2em}
$endif$
$endif$

$if(toc)$
$if(toc-on-own-page)$
% The title block gets a page to itself and the contents start on the next
% one. A contents list that begins halfway down the title page reads as part
% of the title, and the reader cannot tell where the document proper starts.
\newpage
$endif$
\tableofcontents
\newpage
$endif$

$body$

\end{document}
